% !TEX root = ../main.tex
\lecture{8}{2026-11-06}{Ausgabenminimierung, Aggregierte Nachfrage, Produktion und Kosten}
\subsection{Ausgabenminimierung}
Inversion des Problem der Nutzenmaximierung, es wird unter der Bedingung eines spezifischen Nutzen Niveaus, die Ausgaben/Budget minimiert, die diese Bedingung erfüllen.
\begin{itemize}
    \item Exogen gegeben: Pw, Pf, U 
    \item Endogene Variablen: W,F
    \item Gesucht: M
    \item Resultat ist die Hick'sche Nachfrage oder kompensierte Nachfrage
    \begin{itemize}
        \item Es geht darum nach einer Preis Verändung, eine Minimale Veränderung des Budgets zu finden, den das Nutzen niveau erreicht
    \end{itemize}
\end{itemize}

\subsubsection{Grafische Interpretation}
\begin{itemize}
    \item Alle auf der Budget gerade liegenden Kombinationen von W und F sind gleich dem Budget
    \item Alle Parallel verschiebungen der Geraden entsprechen grössere oder kleinere Budgets
    \item Die Indifferenzkurve wird gegeben sein, man muss nur die Iso-Ausgaben-Gerade so verschieben, dass ein Tangentialpunkt entsteht
\end{itemize}

\subsubsection{Formales Optmierungsproblem}
Gegeben eine Nutzensfunktion, Preis zweier Güter und Preisveränderung eines der beiden. Der Optimumspunkt auf der originalen Indifferenzkurve kann mittels Hick'schen Nachfrage berechnet werden oder besser gesagt ist die hick'sche Nachfrage
\begin{itemize}
    \item Die Optimalitätsbedingung für ein gegebenes Nutzenniveau wird nach W oder F gelöst 
    \item Das neue Preis Verhältnis wird eingesetzt
    \item Die neue Budget gerade und entsprechende optimale Kombination der Güter auf der originalen Indifferenzkurve entsteht
    \item Die Differenz der Mengen von W* und F* und originale W und F entspricht den Substitutionseffekt
    \begin{itemize}
        \item Also wie viel vom einen Gut aufgrund der Substitution gegen eines Anderen aufgegeben werden müsste
    \end{itemize}
\end{itemize}

\subsection{Aggregierte Nachfrage}
Bisher nur Nachfrage-Funktionen einzelner Personen berechnet. Um die gesamte Nachfrage Menge in Abhängigkeit vom Preis zu bestimmen müssen die Nachfrage-Funktionen addiert werden.
\begin{itemize}
    \item Horizontale addition 
    \begin{itemize}
        \item Es werden nicht Preise addiert sondern nachgefragten Mengen bei jedem Preis
        \item Nach 0 Punkt, also bei dem Preis, wo keine Person mehr bereit ist zu zahlen, wird Funktion negativ, was konzeptuell nicht gültig ist
        \item Funktion darf nur gelten in positiven Bereichen
        \item Im Fall zweier Funktionen mit unterschiedlicher positiver Bereiche, wird Gesamtnachfrage abschnittweise Definiert 
    \end{itemize}
\end{itemize} 
\section{Produktion und Kosten}

\begin{definition}
    [Produktion] umfasst alle Aktivitäten, die zukünftig Nutzen schaffen. Waren, Dienstleistungen, Intellektuelle/Geistige Güter \footnote{Gar keine Unterscheidung der qualitativ verschiedene Rollen, die die Produktionsbranchen in der Wert Produktion annehmen. Alle Produzieren Gebrauchswerte, aber nicht alle produzieren Wert}
\end{definition}
\textbf{Produktionsfaktoren}
\begin{itemize}
    \item Arbeit
    \item Kapital (Gebäude, Computer) \footnote{Pure Vulgärökonomie. Hier kommt der willkürliche Ausgangspunkt der Ökonomischen Wissenschaft wieder ins Spiel, Kapital ist nicht eine gesellschaftsspezifische Form, die Arbeitsprodukte im produktiven Prozess annehmen, sondern die Objekte sind selbst in ihrer Materie Kapital. Naturalisierung/Mystifizierung gesellschaftlicher Verhältnissen}
    \item Land
    \item Humankapital 
    \item Organisationsstruktur rechtlichen Rahmen
    \item Zwischenprodukte
\end{itemize}

\subsection{Produktionsfunktion}
\subsubsection{Annahmen}
\begin{itemize}
    \item Kapital und Arbeit/Labour (L)
    \item Güter (Q)
    \item Produktionsfunktion F beschreibt Maximale Produktionsmenge in Abhängigkeit der Input-mengen
    \item \begin{equation}
        q = F(L,K)
    \end{equation}
    \item Produktion wird als Blackbox betrachtet, unabhängig der spezifischen Güter, Techniken, usw. \footnote{Also eins der wichtigsten Aspekten der kapitalistischen Produktion wird einfach abstrahiert. Kein einblick in die Wertproduktion, Rolle der Kapitalteile in der Wertreproduktion oder die Rolle der Arbeitskraft in der Mehrwertschöpfung}
\end{itemize}
\subsection{Kurzfristige Produktion}
    Es wird in der kurzen Frist wird Kapital als konstant angenommen, und nur Menge Arbeit variiert. \footnote{In was wird Arbeit und Kapital gemessen? Wie immer in Geld}

\subsubsection{Ertragsgesetzliche Funktion}
Ist eine polynomiale Funktion der Arbeitskraftausgaben.